\documentclass[10pt,a4paper]{article}
\usepackage[margin=2.6cm]{geometry}
\usepackage[T1]{fontenc}
\usepackage{amsmath,amsthm}
\usepackage{newtxtext,newtxmath}
\usepackage{mathpartir,stmaryrd,microtype,booktabs,enumitem,xcolor}
\usepackage{tikz}
\usepackage[colorlinks,linkcolor=black,citecolor=black,urlcolor=black]{hyperref}

\newtheorem{theorem}{Theorem}
\newtheorem{proposition}[theorem]{Proposition}
\theoremstyle{definition}
\newtheorem{definition}[theorem]{Definition}
\newtheorem{example}[theorem]{Example}
\theoremstyle{remark}
\newtheorem{remark}[theorem]{Remark}

\newcommand{\bx}[1]{\Box_{#1}}
\newcommand{\sem}[1]{\llbracket #1\rrbracket}
\newcommand{\sep}{\mathrel{\#}}
\newcommand{\kw}[1]{\mathsf{#1}}
\newcommand{\keep}{\kw{keep}}
\newcommand{\init}{\kw{init}}
\newcommand{\grow}{\kw{grow}}
\newcommand{\any}{\kw{any}}
\newcommand{\val}{\kw{all\text{-}values}}
\newcommand{\st}{\kw{all\text{-}stores}}
\newcommand{\rigid}{\kw{rigid}}
\newcommand{\fresh}{\kw{fresh}}
\newcommand{\old}{\kw{old}}
\newcommand{\dense}{\kw{dense}}
\newcommand{\every}{\kw{every}}
\newcommand{\apart}{\kw{apart}}
\newcommand{\fld}{\kw{all\text{-}fields}}
\newcommand{\cel}{\kw{all\text{-}cells}}
\newcommand{\elt}{\kw{all\text{-}elements}}
\newcommand{\slen}{\kw{strLength}}
\newcommand{\rowt}{\kw{table}}
\newcommand{\esum}{\kw{sum}}
\newcommand{\eff}[1]{\varepsilon_{#1}}
\newcommand{\len}[1]{|#1|}
\newcommand{\tr}[1]{\xrightarrow{#1}}
\newcommand{\resp}{\mathrel{\triangleleft}}
\newcommand{\hoare}[3]{\{#1\}\;#2\;\{#3\}}
\newcommand{\hoarex}[4]{\{#1\}\;#2\;\{#3\mid #4\}}
\newcommand{\code}[1]{\texttt{#1}}
\newcommand{\lean}[1]{\texttt{\small #1}}
\newcommand{\status}[1]{\textsf{\small[#1]}}
\frenchspacing\sloppy

\title{Tolerant frames II\\[2pt]\large What the bound proofs of a text-layout library need}
\author{TScert working note}
\date{13 November 2026}

\begin{document}
\maketitle

\begin{abstract}
\noindent
The text-layout library pretext verifies under TScert only in a rewritten
form. The rewritten form has 422 sites where an index read carries a
check that the original source does not have. This note asks, for each
site, what the proof of its bound needs in the logic of the first paper,
\emph{Tolerant frames}. The 215 sites that need a capability fall under
six capabilities. Two are derivable in the first paper: postconditions
that say by how much a length changed, and a group of rules whose
instances TScert does not derive. A third, finer regions for what a call
may touch, is derivable for two of its three patterns and needs regions
by kind of location for the third. Two need small extensions: bounded
effect variables for a callback passed in, and contracts that hold of a
location or of a type for its whole life. The last needs five richer
term forms. Each extension is a definition with its meaning, one
example, and the statement that makes it sound. One worked program shows
four of the capabilities at marked lines. The analysis's classification
changes for two patterns. Every statement carries one of four marks. \status{paper} names the definition, axiom, rule or theorem of
the first paper that derives it. \status{observed} names the test program
on which TScert's verifier gave the stated verdict. \status{proved on
paper} marks a statement proved in this note. \status{conj.} marks a
conjecture.
\end{abstract}

\section{The question}\label{sec:intro}

The library pretext lays out text: it measures segments of a string into
parallel arrays, and walks those arrays to break lines. Its source
verifies under TScert only after rewrites. A \emph{bound rewrite site} is
an index read \code{x[k]} that the rewritten source writes as
\code{x[k] ?? missing()} or guards by a \code{RangeError}, because no
proof of $0\le k<\len x$ is found there. This note asks which capability
of a logic each such proof lacks. The logic is the one of the first paper,
\emph{Tolerant frames: a modal account of fact preservation without
ownership} (11 November 2026), written TF below; ``TF §$n$'' names its
section $n$. The sites and their minimal programs are those of an analysis
of pretext's rewritten source at commit \code{9268edb1}; every site count below
comes from that analysis.

\paragraph{The count.} Of the 422 bound rewrite sites, 215 are
\emph{blocked}: the proof at the site itself needs a capability. 204 are
\emph{locally supported}: the proof at the site needs nothing new, given
facts established at a builder or a caller. Of these, 39 have the shape of
a verified program \status{observed} and 165 are \status{conj.}. 3 sites
are not placed (Section~\ref{sec:notcovered}). The 215 blocked sites fall
under six capabilities:

\medskip
\noindent
\begin{tabular}{@{}rp{0.36\textwidth}lrp{0.24\textwidth}@{}}
\toprule
 & capability & patterns & sites & in TF\\
\midrule
1 & the effect of a callback passed in & P1 & 50 & no: bounded effect variables\\
2 & posts that say by how much a length changed & P2, P8, P14 & 45 & yes: \textsc{Call}, two-state posts\\
3 & finer regions for what a call may touch & P3, P4, P6 & 49 & P3, P4 yes; P6 no: kind regions\\
4 & facts of a location or a type for good & P5, P15, P16, P20 & 25 & no: contracts\\
5 & richer terms & P9, P10, P12, P17, P19 & 24 & no: five term forms\\
6 & rules TF has and TScert lacks & P7, P11, P13, P18, P21 & 22 & yes\\
\midrule
 & total & & 215 & \\
\bottomrule
\end{tabular}
\medskip

\noindent
A \emph{pattern} P$n$ is the smallest program shape whose bound proof
lacks one capability. The labels P1 to P21 are the analysis's. Each
section below takes one capability. For each pattern it gives the minimal
program, the fact the proof needs in TF's assertion language, where TF
stops, and the smallest extension. Section~\ref{sec:class} gives the
class of each pattern as this note finds it. Section~\ref{sec:sound}
states what makes each extension sound in TF's model.

\section{The notions used}\label{sec:notions}

This section recalls the notions of TF that later sections use, each with
its definition, one example, and its place in TF.

\begin{description}[leftmargin=1.5em,style=sameline]
\item[Heaps and kinds (TF §2)] A heap is a finite map
$h:\mathit{Loc}\rightharpoonup\mathit{Val}\uplus\mathit{Store}$ with
$\mathit{Loc}=\mathbb{N}$ and $\mathrm{dom}\,h=\{0,\dots,\len h-1\}$.
Its domain is partitioned by kind: $\mathrm{dom}\,h=F\uplus C\uplus
E\uplus S$, the field locations, the cells, the element cells and the
store locations. The first three hold values and form the \emph{value}
class; $S$ holds stores and forms the \emph{store} class. A cell is the
location of a module binding, or of a local binding that a closure
writes or sees change; every other local binding is a variable of the
environment $\rho:\mathit{Var}\rightharpoonup\mathit{Val}$. Example: after
\code{const o = \{ n: 1 \};} the object \code{o} is the map
$\{\code{n}\mapsto l\}$ for one $l\in F$, and \code{o} is a variable.
\item[Stores and lengths (TF §2)] $\mathit{Store}=(\mathit{Loc}\uplus\{\circ\})^{*}$.
For $\sigma\in\mathit{Store}$, $\len\sigma$ is its length and
$\mathit{cells}(\sigma)=\{c_i\mid c_i\ne\circ\}$. An array is an object
with one distinguished field whose location $s\in S$ holds its store.
Example: \code{const a = [10, , 30];} gives $h(s)=\langle
l_0,\circ,l_2\rangle$, length 3, $l_0,l_2\in E$.
\item[Contracts and immutable locations (TF §2, §3)] A world $W$ maps each
location to its \emph{contract}, the set of contents it may hold. A
location is immutable when its contract admits one value. TF has two
kinds: a field declared \code{@stable}, and a typed array's length,
written $\kw{view}(w)$. Example: in
\code{type Runs = \{ /** @stable */ readonly openBrackets: number[] \}}
the field \code{openBrackets} of each \code{Runs} object always holds the
same array.
\item[Admissible transitions (TF §3)] $h\tr{}_\rho h'$ when $h'$ is the
heap of a state whose world extends that of $h$: every location keeps its
contract, and $h'$ meets it. Every run of typed code is admissible
\status{paper: TF §3, \emph{Admissible transitions}}. Its properties are
(T1) preorder, (T2) growth $\len h\le\len{h'}$, (T3) fixed objects: an
object's keys and field locations do not depend on the heap, and (T4)
immutable locations: a location whose contract admits the single value
$u$ holds $u$ in $h$ and $h'$. Example: \code{o.n = 2} on the object
above is admissible when the contract of $l$ is ``a number''.
\item[Changes (TF §3)] A change $\delta$ is a family of preorders
$\delta_b$ on contents; $\keep\sqsubseteq\init\sqsubseteq\kw{scalar}\sqsubseteq\grow\sqsubseteq\any$.
$\grow_b$ relates two stores when the second has a length at least the
first's, keeps every dense prefix, and lists no old location the first
does not list. Example: $\grow_6$ relates $\langle4,\circ,5\rangle$ to
$\langle4,\circ,5,6\rangle$ and not to $\langle4\rangle$.
\item[Regions (TF §3)] A path $\pi=x.f_1\ldots f_k$ is a variable
followed by field names. A region $r$ denotes a set $r(\rho,h)$ of
locations: $\st$ denotes $S$; $\val$ denotes $F\cup C\cup E$;
$\pi.\kw{store}$ denotes the store location of the array $\pi$ reads;
$\pi[\ast]$ denotes it with the cells its store lists; $\pi.f$ denotes one
field location. Example: for \code{a} above, $a[\ast]$ denotes
$\{s,l_0,l_2\}$.
\item[Effects and their action (TF §3)] An effect $\Phi$ is a finite set
of entries $r:\delta$. $h\tr{\Phi}_\rho h'$ when $h\tr{}_\rho h'$ and
every location $l<\len h$ holds what it held or lies in a region
$r(\rho,h)$ of an entry $r:\delta$ and changed within $\delta_{\len h}$.
$\top=\{\val:\any,\ \st:\any\}$. $\Phi$ is \emph{rigid} when its action is
transitive; effects over $\val$ and $\st$ are rigid. Example: the effect
of \code{out[j] = v} is $\{\val:\any,\ \st:\grow\}$.
\item[Assertions (TF §4)] Atoms are linear inequalities over variables,
lengths $\len\pi$ and $\kw{view}(\pi)$; equalities $\pi\doteq v$;
$\dense(x,n)$; kinds of contents; and $\every(z,t,\varphi)$, each
element of $z$ below index $t$ satisfies $\varphi$, a formula of that one
element. Assertions add $\wedge$, $g\vee F$ with $g$ invariant,
$\exists$, $\pi\sep\pi'$ (different store locations), $\apart(x,y)$ (no
common element cell), $\fresh(x)$ and the modality. Example:
$\every(a,1,v>0)$ holds of \code{[10, , 30]}.
\item[The modality (TF §4)] $\rho,m,h\models\bx{\Phi}F$ when $F$ holds in
every $h'$ with $h\tr{\Phi}_\rho h'$. $F$ is \emph{invariant} when
$F\vdash\bx{\top}F$; $\rigid(\pi)\triangleq\exists v.\,\bx{\top}(\pi\doteq
v)$. Example: $3\le\len a$ entails $\bx{\val:\any,\ \st:\grow}(3\le\len
a)$, and does not entail $\bx{a.\kw{store}:\any}(3\le\len a)$.
\item[Support, respect, tolerant frame (TF §5)] A support $S$ of $F$ is a
set of pairs $r:\delta$ such that $F$ survives every admissible
transition that changes each location of each $r$ within its $\delta$.
$\Phi\resp_K S$ when each entry $r':\gamma$ of $\Phi$ and each pair
$r:\delta$ of $S$ have $\gamma\sqsubseteq\delta$ or regions disjoint
under $K$. Then $K\wedge F\vdash\bx{\Phi}F$ \status{paper: Theorem
\emph{Tolerant frame}}. Example: $\{x.\kw{store}:\grow\}$ supports
$n\le\len x$.
\item[Axioms (TF §5)] \textsc{(Grow)} keeps a lower bound of a length
under $\{\val:\any,\ \st:\grow\}$; \textsc{(Dense)} a dense prefix under
the same effect; \textsc{(Sep)} a fact of $\len x$ under
$\{\val:\any,\ y.\kw{store}:\any\}$ given $x\sep y$; \textsc{(Every)} an
$\every$ over $z$ under $\{x[\ast]:\any\}$ given $\apart(x,z)$;
\textsc{(Init)} a length under $\{\val:\init\}$; \textsc{(Rigid)} a fact
of a variable's array for the array at a rigid path; \textsc{(View)} a
fact of $\kw{view}(w)$ under $\top$. Example: \textsc{(Grow)} gives
$n\le\len x\vdash\bx{\val:\any,\ \st:\grow}(n\le\len x)$.
\item[Triples and rules (TF §6)] $\hoarex PcQ{Q^e}$ is partial
correctness; $Q$ is two-state, with $\old(e)$ the value of $e$ where $c$
started. The rules are \textsc{Frame}, \textsc{Frame}$^{+}$,
\textsc{Loop}, \textsc{Call}, \textsc{Alloc}, \textsc{Seq-Eff},
\textsc{Fresh-Eff}. Example: \code{x.push(v)} has the post
$\len x=\old(\len x)+1$.
\item[Checked transfer (TF §7)] An edge $e$ of a control-flow graph
carries an effect $\Phi_e$, a sound two-state post $Q_e$, a condition and
a renaming. A fact crosses $e$ when it is derivable under
$\bx{\Phi_e}$; the rest is derived from $Q_e$. Example: the node $j'=j+1$
has the post $j'=j+1$.
\end{description}

\noindent
The sets and maps of the account are these.

\begin{definition}[Sites]\label{def:sites}
$R$ is the set of bound rewrite sites, $\len R=422$. It is partitioned as
$R=B\uplus L\uplus U$: the blocked sites $B$, $\len B=215$; the locally
supported sites $L$, $\len L=204$; the sites not placed $U$, $\len U=3$.
$\mathcal P=\{\mathrm P1,\ldots,\mathrm P21\}$ is the set of patterns.
The map $p:B\to\mathcal P$ gives a blocked site its first blocker, the
first capability its own read lacks. The map
$\kappa:\mathcal P\to\{1,\ldots,6\}$ gives a pattern its capability.
\end{definition}
\noindent
Example: the read \code{starts[i + 1]} of \code{measureAnalysis} at line
378 of \code{prepare.ts} is a site $r\in B$ with $p(r)=\mathrm P4$ and
$\kappa(\mathrm P4)=3$.

\begin{definition}[Classes]\label{def:classes}
Let $F$ be the fact a pattern's bound proof needs, and $c$ the step it
must cross. The pattern is in one class:
\begin{description}[leftmargin=1.5em,style=sameline]
\item[(a) not stated] No assertion of TF states $F$. Example:
$\len{\mathit{rows}}=\mathit{states}\cdot\mathit{width}$, a product of two
variables.
\item[(b) not kept] $F$ is an assertion, and for no effect $\Phi$ that TF
can state for $c$ is $\bx{\Phi}F$ derivable. Example: an $\every$ over a
\code{number[]} across a call that writes a module cell, whose only TF
effect is $\{\val:\any\}$.
\item[(c) kept in TF] $F$ is an assertion and its proof across $c$ is
derivable in TF, from posts sound for the steps involved. Example: a
length across a call whose two-state post states the length.
\item[(d) a contract] $F$ holds because of the contract of a location or
an invariant of a type, beyond the single-value contracts of TF.
Example: the length of a module constant array literal.
\end{description}
\end{definition}

\noindent
Two notions beyond TF are used throughout.

\begin{definition}[String length]\label{def:slen}
For a variable $x$ that holds a string, $\slen(x)\in\mathbb{N}$ is the
number of UTF-16 code units of $\rho(x)$. Linear terms may contain
$\slen(x)$.
\end{definition}
\noindent
A string is a scalar, so $\slen(x)$ reads no location: its support is
$\emptyset$, and every atom over variables and string lengths is
invariant. Example: after \code{const s = "ab"}, $\slen(s)=2$ holds after
every step. TF's atoms have array lengths and $\kw{view}$ only. TScert
has this term \status{observed: \code{t08\_strings.ts}, \code{goodCode}
verified}.

\begin{definition}[Paths from cells]\label{def:cellpath}
A path may start at a binding $b$ modelled as a cell $l_b\in C$. Its read
starts from $h(l_b)$.
\end{definition}
\noindent
Example: for the module binding \code{let bmp = null}, the path
$\mathit{bmp}$ reads $h(l_{\mathit{bmp}})$. TF uses this reading for the
cell \code{c} of its write-once cache pattern (TF §8) and does not define
it. A fact through such a path reads $l_b$, so it survives an effect
disjoint from $l_b$, or any effect when a contract fixes $l_b$
(Definitions~\ref{def:imm} and~\ref{def:cellinv}).

\section{A worked program}\label{sec:worked}

The program below builds two parallel arrays and walks them. It is this
note's own, in the shape of pretext's builder and walker. A mark $(n)$
names the capability whose proof the line needs.

\medskip
\noindent
{\small
\begin{tabular}{@{}r@{\hspace{1em}}l@{\hspace{1em}}l@{}}
1 & \code{type Lists = \{ /** @stable */ readonly widths: number[];} & \\
2 & \code{~~~~~~~~~~~~~~~/** @stable */ readonly kinds: number[] \};} & \\
3 & \code{const KINDS: readonly string[] = ["text", "space", "break"];} & \\
4 & \code{let calls: number = 0;} & \\
5 & \code{function measure(text: string): number \{ calls++; return text.length; \}} & \\
6 & \code{function append(lists: Lists, texts: string[]): void \{} & \\
7 & \code{~~for (let i = 0; i < texts.length; i++) \{} & \\
8 & \code{~~~~const w = measure(texts[i]!);} & (3)\\
9 & \code{~~~~lists.widths.push(w);} & \\
10 & \code{~~~~lists.kinds.push(texts[i] === " " ? 1 : w === 0 ? 2 : 0);} & (6)\\
11 & \code{~~\}} & \\
12 & \code{\}} & \\
13 & \code{function build(texts: string[]): Lists \{} & \\
14 & \code{~~const lists: Lists = \{ widths: [], kinds: [] \};} & \\
15 & \code{~~append(lists, texts);} & (2)\\
16 & \code{~~return lists;} & \\
17 & \code{\}} & \\
18 & \code{function walk(lists: Lists, visit?: (i: number) => void): number \{} & \\
19 & \code{~~const n = lists.widths.length; let x = 0;} & \\
20 & \code{~~for (let i = 0; i < n; i++) \{} & \\
21 & \code{~~~~x = lists.kinds[i]! === 2 ? 0 : x + lists.widths[i]!;} & \\
22 & \code{~~~~visit?.(i);} & (1)\\
23 & \code{~~\}} & \\
24 & \code{~~return x;} & \\
25 & \code{\}} & \\
26 & \code{export function lastKind(texts: string[]): string \{} & \\
27 & \code{~~const lists = build(texts);} & \\
28 & \code{~~const x = walk(lists, undefined);} & (1)\\
29 & \code{~~const n = lists.kinds.length;} & \\
30 & \code{~~return n === 0 ? "" : KINDS[lists.kinds[n - 1]!]! + x;} & (4)\\
\end{tabular}}
\medskip

\paragraph{The heap.} Number the locations in allocation order, the
module first. At line 28 of \code{lastKind(["ab", " "])}, before the call
of \code{walk}, one heap $h_{28}$ the program can have is:
\[
\begin{array}{llll}
\text{location} & \text{kind} & \text{holds} & \\
0 & C & k & \text{the cell of \code{KINDS}}\\
1 & S & \langle 2,3,4\rangle & \text{the store of \code{KINDS}}\\
2,3,4 & E & \code{"text"},\ \code{"space"},\ \code{"break"} & \text{its elements}\\
5 & C & 2 & \text{the cell of \code{calls}}\\
6 & S & \langle 7,8\rangle & \text{the store of \code{texts}}\\
7,8 & E & \code{"ab"},\ \code{" "} & \text{its elements}\\
9 & S & \langle 13,15\rangle & \text{the store of \code{lists.widths}}\\
10 & S & \langle 14,16\rangle & \text{the store of \code{lists.kinds}}\\
11 & F & W & \text{the field \code{lists.widths}}\\
12 & F & K & \text{the field \code{lists.kinds}}\\
13,15 & E & 2,\ 1 & \text{\code{lists.widths[0]}, \code{lists.widths[1]}}\\
14,16 & E & 0,\ 1 & \text{\code{lists.kinds[0]}, \code{lists.kinds[1]}}
\end{array}
\]
Here $k$, $W$ and $K$ are the arrays whose store locations are 1, 9
and 10. $C=\{0,5\}$, $S=\{1,6,9,10\}$, $F=\{11,12\}$, and $E$ is the rest.
The environment holds \code{texts} and
$\mathit{lists}=\{\code{widths}\mapsto 11,\ \code{kinds}\mapsto 12\}$.
The fields 11 and 12 are \code{@stable}, so $\rigid(\mathit{lists}.\mathit{widths})$
and $\rigid(\mathit{lists}.\mathit{kinds})$ hold \status{paper: TF §4,
\emph{Invariant assertions and rigid paths}}.

\paragraph{The facts.} The read at line 30 needs, at line 28, the three
facts
\[
\begin{array}{ll}
F_1 = \every(\mathit{lists}.\mathit{kinds},\ \len{\mathit{lists}.\mathit{kinds}},\ 0\le v<3) & \text{the kinds index \code{KINDS}},\\
F_2 = (\len{\mathit{KINDS}}=3) & \text{the length of the module table},\\
F_3 = (\len{\mathit{lists}.\mathit{widths}}=\len{\mathit{lists}.\mathit{kinds}}) & \text{the arrays are parallel}.
\end{array}
\]
$F_3$ is \code{walk}'s precondition, and $F_1$ and $F_2$ must survive the
call at line 28. Each later section returns to its marked line:
line 22 and line 28 in Section~\ref{sec:cap1}, line 15 in
Section~\ref{sec:cap2}, line 8 in Section~\ref{sec:cap3}, line 30 in
Section~\ref{sec:cap4}, and line 10 in Section~\ref{sec:cap6}.

\paragraph{The verdicts.} As printed, with no annotation beyond the
types, TScert verifies \code{measure}, \code{append} and \code{build}:
they push and read \code{texts[i]} below its length, and no read of
theirs needs a fact from elsewhere. It refuses \code{walk} at the read
\code{lists.kinds[i]} of line 21, which needs $F_3$, a precondition the
listing does not write. It refuses \code{lastKind} at the read
\code{KINDS[lists.kinds[n - 1]!]} of line 30, which needs $F_1$ and
$F_2$ \status{observed: branch \code{single-pass} at \code{b7c8a559},
13 November 2026}. Line 4 writes the type of \code{calls} because TScert
requires a module \code{let} to state its type.

\section{The effect of a callback passed in}\label{sec:cap1}

\subsection{P1: a visitor of unknown effect between reads (50 sites)}

\emph{Program.} \code{function walk(xs: number[], visit: Visitor |
undefined): number} reads \code{xs[i]} in a loop and calls
\code{visit(s)} when \code{visit !== undefined}. The function \code{bad}
calls \code{walk(xs, undefined)} and then reads \code{xs[i]} under the
precondition $0\le i<\len{\mathit{xs}}$ \status{observed:
\code{t02\_optional\_callback.ts}}.

\emph{Fact.} After the call, $i<\len{\mathit{xs}}$. In pretext's walker
\code{walkPreparedComplexLines}, with $c$ the segment count read at
entry, each loop head needs
$c\le\len{\mathit{widths}}\wedge 0\le\mathit{cursor}.\mathit{segmentIndex}<c$
across the call \code{onLine?.(\ldots)}.

\emph{Where TF stops.} The parameter's type states no effect, so the
effect of a call of \code{visit} is $\top$, and so is the effect of
\code{walk}. The transition that sets \code{xs.length = 0} is bounded by
$\top$ and breaks the fact. TF states one closed bound per callback
parameter: an effect has no variables (TF §10, \emph{Effects are
closed}). A closed bound keeps lower bounds of lengths: with the bound
$\Psi=\{\val:\any,\ \st:\grow\}$, $i<\len{\mathit{xs}}$ survives
\status{paper: \textsc{(Grow)}} \status{observed:
\code{t02\_optional\_callback.ts}, \code{goodBounded} verified}. Two
things remain out of reach. Inside the walker, the fact
$0\le\mathit{cursor}.\mathit{segmentIndex}$ reads a field location; every
closed bound that admits the application's visitors contains
$\val:\any$, which covers the field, and no TF region excludes one
object's fields from $\val$. At a caller passing no visitor, the walker
calls nothing, yet under a closed bound its effect is still $\Psi$, so
every fact over value locations and every upper bound of a length is lost
at the call. The first needs a region; the second needs an effect
variable.

\begin{definition}[Complement and object regions]\label{def:compl}
For a path $\pi$ that reads an object $v$, the region $\pi.\ast$ denotes
$\{v(f)\mid f\in\mathrm{dom}\,v\}$, the field locations of $v$. For
regions $r_1,r_2$, the region $r_1\setminus r_2$ denotes
$r_1(\rho,h)\setminus r_2(\rho,h)$.
\end{definition}
\noindent
Example: if $\mathit{cursor}=\{\code{segmentIndex}\mapsto 20,\ \code{graphemeIndex}\mapsto 21\}$,
then $\mathit{cursor}.\ast$ denotes $\{20,21\}$, and
$\val\setminus\mathit{cursor}.\ast$ denotes every value location but 20
and 21. Under the bound
$\Psi_{\mathit{onLine}}=\{\val\setminus\mathit{cursor}.\ast:\any,\ \st:\grow\}$
the fact $0\le\mathit{cursor}.\mathit{segmentIndex}<c$, of support
$\{\mathit{cursor}.\mathit{segmentIndex}:\keep\}$, survives by the
Tolerant frame theorem: one entry is disjoint from the field, and the
other has the store class.

\begin{definition}[Bounded effect variable]\label{def:effvar}
Let $p$ be a parameter of function type of a function $f$. An
\emph{effect variable} $\eff p$ stands for the effect of the value $p$
holds. A signature of $f$ has the form
\[
\forall\eff p\sqsubseteq\Psi.\quad \hoare{P}{f(\vec p)}{Q} : \Phi\cup\eff p ,
\]
where the \emph{bound} $\Psi$ is a closed rigid effect over class
regions and regions at rigid paths of $f$'s parameters. Inside $f$ a call of $p$ has the effect
$\eff p$. At a call of $f$ the variable is instantiated: to $\emptyset$
when the argument is \code{undefined}, and to $\Psi'$ when the argument
is a function of effect $\Psi'$, with the obligation
$\Psi'\sqsubseteq\Psi$. The effect of the call is
$\Phi\cup\Psi'[\vec x/\vec p]$.
\end{definition}

\noindent
The rule it gives at a call, \textsc{Call-Var}, is \textsc{Call} at the instance:
\begin{mathpar}
\inferrule*[left=Call-Var]{f(\vec p):\forall\eff p\sqsubseteq\Psi.\ \hoare PfQ:\Phi\cup\eff p \\ x_p\text{ holds a function of effect }\Psi'\text{, or \code{undefined} and }\Psi'=\emptyset \\ \Psi'\sqsubseteq\Psi}{\hoare{P[\vec x/\vec p]\wedge\bx{(\Phi\cup\Psi')[\vec x/\vec p]}R}{y=f(\vec x)}{Q[\vec x/\vec p]\wedge R}}
\end{mathpar}

\noindent
Inside $f$, a fact survives a call of $p$ for every admitted $\eff p$
exactly when it survives $\Psi$, by law (c) of TF's Proposition
\emph{Laws} and because $\Psi$ itself is admitted. So the variable adds
nothing to the walker's own proof; it makes the caller's effect exact.

\emph{The worked program.} At line 22,
$\code{walk}:\forall\eff{\mathit{visit}}\sqsubseteq\{\val:\any,\ \st:\grow\}.\ \emptyset\cup\eff{\mathit{visit}}$.
The loop facts $n\le\len{\mathit{lists}.\mathit{widths}}$ and
$n\le\len{\mathit{lists}.\mathit{kinds}}$ are lower bounds at rigid paths
and survive $\Psi$ \status{paper: \textsc{(Grow)}, \textsc{(Rigid)},
\textsc{Loop}}. At line 28 the argument is \code{undefined}, so the call
has the effect $\emptyset$, and $F_1$ and $F_2$ hold after it by
\textsc{Frame} \status{proved on paper: Section~\ref{sec:sound}}.
Without the variable, the effect at line 28 is $\Psi$. The visitor
\code{i => \{ lists.kinds[0] = 7; \}} is bounded by $\Psi$, it changes
location 14 of $h_{28}$ to 7, and line 30 then reads
\code{KINDS[7]}, which is \code{undefined}.

\emph{Sites.} 47 sites lie inside the two walkers that call
\code{onLine}: \code{walkPreparedLinesRaw} (310) and
\code{walkPreparedComplexLines} (46 sites, 584 to 1085). They need a
closed bound with Definition~\ref{def:compl}. 3 sites lie at callers that
pass no visitor: \code{countSteppedLines} (407) and
\code{prepareRichInline} (368, 496). They need
Definition~\ref{def:effvar}.

\emph{Evidence.} \code{bad} is refused at \code{xs[i]}, and
\code{goodBounded} and \code{goodPlain} are verified \status{observed:
\code{t02\_optional\_callback.ts}}. \code{flagsWalk}, which reads a
\code{Uint8Array} below a count taken at entry across a visitor of
unknown effect, is verified by \textsc{(View)}; \code{widthsWalk}, the
same loop over a \code{number[]}, is refused \status{observed:
\code{t13\_view\_callback.ts}}.

\begin{remark}[Classification of P1]
The analysis gives P1 one definition. This note finds two needs. The 47
walker sites are class (b) for want of Definition~\ref{def:compl}: the
fact is stated, and TF can state no bound that admits the application's
visitors and keeps the cursor's fields. The 3 caller sites are class (b)
for want of Definition~\ref{def:effvar}. The bound on a visitor that the
application passes is a trust decision in both (Section~\ref{sec:notcovered}).
\end{remark}

\section{Posts that say by how much a length changed}\label{sec:cap2}

These patterns are in TF. A step's post may be two-state
(TF §6, \emph{Triples}), and \textsc{Call} carries the callee's post
$Q[\vec x/\vec p]$ to the caller \status{paper: \textsc{Call}}. An exact
post constrains the written location only; every other fact crosses the
step by \textsc{Frame} (TF §6, \emph{Exact posts}).

\subsection{P2: growth by a known count across a callee (31 sites)}

\emph{Program.} \code{measure(lists, n)} pushes $n$ values onto
\code{lists.widths}, a \code{@stable} field, with the effect
$\{\val:\any,\ \mathit{lists}.\mathit{widths}[\ast]:\any\}$. The caller
calls \code{measure(lists, 10)} in a loop on \code{j += 10} with
\code{j + 10 <= cap}, and reads \code{flags} at the last index of
\code{lists.widths}, where \code{flags = new Uint8Array(cap)}.
The program as TScert verifies it \status{observed: the case
\code{bounds-pattern-known-growth-refused}, verified on the branch
\code{single-pass}, 13 November 2026}:

\medskip
\noindent
{\small
\begin{tabular}{@{}r@{\hspace{1em}}l@{}}
1 & \code{type Lists = \{ /** @stable */ readonly widths: number[] \};} \\
2 & \code{/** @requires safeInteger(n) \&\& 0 <= n \&\& n <= 1000 \&\& 0 <= lists.widths.length} \\
3 & \code{~*~ @ensures lists.widths.length === old(lists.widths.length) + n */} \\
4 & \code{function measure(lists: Lists, n: number):} \\
5 & \code{~~~~Modifies<\{ values: "any"; stores: "keep"; "lists.widths[*]": "any" \}, void> \{} \\
6 & \code{~~/** @invariant safeInteger(k) \&\& safeInteger(n) \&\& 0 <= k \&\& k <= n} \\
7 & \code{~~~*~~~~~~~~~~~~\&\& lists.widths.length === old(lists.widths.length) + k */} \\
8 & \code{~~for (let k = 0; k < n; k++) lists.widths.push(k);} \\
9 & \code{\}} \\
10 & \code{/** @requires safeInteger(cap) \&\& 0 <= cap \&\& cap <= 100000 */} \\
11 & \code{export function grown(cap: number): number \{} \\
12 & \code{~~const widths: number[] = [];} \\
13 & \code{~~const lists: Lists = \{ widths \};} \\
14 & \code{~~const flags = new Uint8Array(cap);} \\
15 & \code{~~let s = 0;} \\
16 & \code{~~/** @invariant safeInteger(j) \&\& 0 <= j \&\& safeInteger(cap)} \\
17 & \code{~~~*~~~~~~~~~~~~\&\& 0 <= lists.widths.length \&\& lists.widths.length <= j \&\& j <= cap */} \\
18 & \code{~~for (let j = 0; j + 10 <= cap; j += 10) \{} \\
19 & \code{~~~~measure(lists, 10);} \\
20 & \code{~~~~s += flags[lists.widths.length - 1]!;} \\
21 & \code{~~\}} \\
22 & \code{~~return s;} \\
23 & \code{\}} \\
\end{tabular}}
\medskip

\noindent
Line 3 is the two-state postcondition: \code{old(e)} is the value of
\code{e} where \code{measure} was entered. The read under proof is at
line 20. Each other annotation is needed, for these reasons.
\begin{description}[leftmargin=1.5em,style=sameline]
\item[Line 2, \code{n} a nonnegative safe integer] Without it the
postcondition is false: \code{measure(lists, 1.5)} pushes two values, and
a negative \code{n} pushes none.
\item[Line 2, \code{0 <= lists.widths.length}] It states that the length
is defined at entry. In TScert's model this does not follow from the type
\code{Lists}: a value meets its type only to the depth of the step
budget \status{observed: \code{measure} is refused at its entry without
it}.
\item[Line 5, the effect] The caller keeps its facts about \code{flags}
across the call because the call changes no store but
\code{lists.widths}.
\item[Lines 6--7 and 16--17, the loop invariants] Rule \textsc{Loop}
needs an invariant at each loop head.
\item[Line 20, the read through \code{lists.widths}] With the read written
\code{flags[widths.length - 1]}, through the local name, \code{grown} is
refused at the read \status{observed}. TF derives that form too, from
$\mathit{lists}.\mathit{widths}\doteq\mathit{widths}$ and
$\rigid(\mathit{lists}.\mathit{widths})$; why TScert does not is not
established.
\end{description}

\emph{Fact.} The loop invariant $\len{\mathit{widths}}=j$, with
\code{lists.widths} the array \code{widths}. It needs the callee's post
\[
\hoare{\kw{true}}{\code{measure}(\mathit{lists},n)}{\len{\mathit{lists}.\mathit{widths}}=\old(\len{\mathit{lists}.\mathit{widths}})+n} .
\]

\emph{In TF.} \textsc{Call} gives
$\len{\mathit{lists}.\mathit{widths}}=j+10$ after the call;
$\rigid(\mathit{lists}.\mathit{widths})$ and
$\mathit{lists}.\mathit{widths}\doteq\mathit{widths}$ give
$\len{\mathit{widths}}=j+10\le\mathit{cap}=\kw{view}(\mathit{flags})$,
and \textsc{(View)} keeps $\kw{view}(\mathit{flags})=\mathit{cap}$. The
read at $j+10-1$ is in bounds \status{paper: \textsc{Call},
\textsc{(Rigid)}, \textsc{(View)}, \textsc{Loop}}. Class (c).

\emph{The worked program.} At line 15, \code{append} has the post
\[
\len{\mathit{lists}.\mathit{widths}}=\old(\len{\mathit{lists}.\mathit{widths}})+\len{\mathit{texts}}
\ \wedge\
\len{\mathit{lists}.\mathit{kinds}}=\old(\len{\mathit{lists}.\mathit{kinds}})+\len{\mathit{texts}} .
\]
Both arrays are empty at line 14 \status{paper: \textsc{Alloc}}, so $F_3$
holds after line 15. A one-state post of \code{append} states lower
bounds only, and $F_3$ does not follow from them.

\emph{Sites.} 30 in \code{prepareRichInline}, where each read or guarded
write at \code{at = widths.length} needs
$\len{\mathit{widths}}\le\kw{view}(\mathit{flags})$ after
\code{measureAnalysis} appended \code{to - from} segments; 1 in
\code{measureAnalysis} (631).

\emph{Evidence.} The program above is verified. With the callee growing
the list by more than the loop's step, the read is refused, and the
program is unsafe under Node \status{observed: the case
\code{invalid-bounds-pattern-known-growth}}. With the pushes written in
the caller's own body, the program is verified, each \code{push} having
the post $\len{\mathit{widths}}=\old(\len{\mathit{widths}})+1$
\status{observed: the case \code{bounds-pattern-known-growth}}.

\subsection{P8: the length after a write at the length (10 sites)}

\emph{Program.} \code{xs[k] = 1; return xs[k]!;} under
$0\le k\le\len{\mathit{xs}}$ \status{observed: \code{t17\_write\_len.ts}}.

\emph{Fact.} After the write, $k<\len{\mathit{xs}}$ and
$\dense(\mathit{xs},k+1)$ when $\dense(\mathit{xs},k)$ held before.

\begin{proposition}[Write at the length]\label{prop:writelen}\status{proved on paper}
For an array $x$ and $0\le k\le\len x<N$,
\[
\hoare{\kw{true}}{x[k]=v}{\len x=\max(\old(\len x),k+1)\ \wedge\ (\old(\dense(x,k))\Rightarrow\dense(x,k+1))} .
\]
\end{proposition}
\begin{proof}
By TF's table of writes (TF §2). If index $k$ is present, the write
changes its element cell and leaves the store. If $k$ is a hole below the
length, the store is filled at $k$ and keeps its length. If $k=\len x$,
the store gains one new cell at index $k$. In each case index $k$ is
present after the write, every other index keeps its presence, and the
length is $\max(\old(\len x),k+1)$.
\end{proof}

\noindent
In TF this post is a sound two-state post of the step, which checked
transfer admits (TF §7). TF does not state it. Class (c). A builder that
appends by index, \code{widths[at] = w} with \code{at} the current
length, keeps $\mathit{at}\le\len{\mathit{widths}}$ at its loop head with
it.

\emph{Sites.} 9 in \code{measureAnalysis} (\code{prepare.ts}) and 1 in
\code{prepareRichInline} (646).

\emph{Evidence.} \code{readWritten}, \code{readBefore} and
\code{readBelow} are verified \status{observed:
\code{t17\_write\_len.ts}, branch \code{single-pass} at \code{b7c8a559}}.
\code{append} and \code{appendNoCall}, builders that append by
\code{widths[at] = mi}, are verified \status{observed:
\code{t16\_supported.ts}, same branch}. TScert derives this post.

\begin{remark}[The post of P8]
The analysis writes this post with the conjunct $x[k]\doteq v$. That
conjunct reads an element at an index, which TF's paths do not express
(Definition~\ref{def:ipath} adds it). The bound proofs of the 10 sites
need only the length and the density, which TF states.
\end{remark}

\subsection{P14: length contracts of string and array natives (4 sites)}

\emph{Program.} \code{const normalized = text.replace(/ +/g, " ");}
followed by reads of \code{breaks = new Uint8Array(text.length + 1)} at
indices below \code{normalized.length} \status{observed:
\code{t08\_strings.ts}}.

\emph{Fact.} $\slen(\mathit{normalized})\le\slen(\mathit{text})$, the post
of the native step. With Definition~\ref{def:slen} it is invariant once
stated. What is missing is the post itself. The posts the sites need
are:
\[
\begin{array}{ll}
\code{r = s.replace(e, t)},\ \text{each match no shorter than }t & \slen(r)\le\slen(s)\\
\code{r = s.split(d)} & 1\le\len r\\
\code{r = xs.map(f)} & \len r=\old(\len{\mathit{xs}})\\
\code{r = w.subarray(b, e)},\ 0\le b\le e\le\kw{view}(w) & \kw{view}(r)=e-b
\end{array}
\]
The post of \code{map} holds whatever \code{f} does: \code{map} reads the
length once at its start (TF §6, Proposition \emph{Reliance}). Each post
is an assumption about a native, a trust decision
(Section~\ref{sec:notcovered}). Class (c).

\emph{Sites.} \code{getGeckoLineBreaks} (684), \code{getWebKitLineRules}
(784, a \code{split}), and \code{getFontData} (145, 146, a \code{map} of
a ten-element literal). 33 locally supported sites depend on the
\code{replace} and \code{slice} posts of \code{analyzeText}.

\emph{Evidence.} \code{bad} is refused at its loop's entry edge;
\code{goodSlice}, a \code{slice} with an exact length contract, is
verified \status{observed: \code{t08\_strings.ts}}.

\begin{remark}[String lengths in P11 and P14]
The facts of P14's minimal program, of site 684 and of P11 are about
string lengths, which TF's atoms do not express. They are class (c) with
Definition~\ref{def:slen}, a term of empty support. The other three sites
of P14 are about array lengths and are class (c) in TF.
\end{remark}

\section{Finer regions for what a call may touch}\label{sec:cap3}

A fact survives a step when the step's effect respects a support of the
fact. Two of the three patterns here are instances of the Tolerant frame
theorem; one needs a new region.

\subsection{P3: a quantified fact across a write outside its array (19 sites)}

\emph{Program.} A loop over \code{items: Int32Array} writes
\code{fragment.width = s} of a local object and reads
\code{widths[items[i]!]} \status{observed:
\code{t03\_every\_fresh\_field.ts}}.

\emph{Fact.} $F=\every(\mathit{items},\kw{view}(\mathit{items}),0\le
v<\len{\mathit{widths}})$ across the write, whose effect is
$\{\mathit{fragment}.\mathit{width}:\any\}$.

\emph{In TF.} $S=\{\mathit{items}[\ast]:\keep,\ \mathit{widths}.\kw{store}:\grow\}$
supports $F$: the fact reads the store and the cells of \code{items}, and
a lower bound of the length of \code{widths}; $\kw{view}(\mathit{items})$
is immutable. The region $\mathit{fragment}.\mathit{width}$ denotes a
location of $F$-kind, and the two regions of $S$ denote locations of
$S\cup E$. They are disjoint under $K=\kw{true}$, so
$F\vdash\bx{\mathit{fragment}.\mathit{width}:\any}F$ \status{paper:
Theorem \emph{Tolerant frame}}. Class (c).

\emph{Sites.} 12 in \code{createLine} (\code{rich-inline.ts}), between a
field write of \code{fragment} and \code{fragments.push(fragment)}; 7 in
\code{setEmptyObjectFacts}, across element writes of \code{spaces}, a
\code{number[]}, and \code{returns}, a \code{Uint8Array}. For the 7, the
effect is $\{\mathit{spaces}[\ast]:\any\}$, and respect needs
$K=\apart(\mathit{spaces},z)$ for each array $z$ the facts read. TF gives
$\apart$ by \textsc{Alloc} at the builder, where the arrays are created;
the facts must then carry it to these sites \status{conj.}. Different
element types give no $\apart$ in TF.

\emph{Evidence.} \code{bad} is refused at \code{widths[items[i]!]}, also
with the fact as a loop invariant; \code{good}, the loop without the
write, is verified \status{observed: \code{t03\_every\_fresh\_field.ts},
\code{t03b\_every\_inv.ts}}.

\begin{remark}[Two sites of P3]
The read at line 1185 of \code{createLine} (2 sites) also needs
$\every(\mathit{fragments},k,\mathit{fragments}[i].\mathit{itemIndex}<n)$,
a fact about a field of each element. TF lists that form outside its
language (TF §8). So these 2 sites are class (a) in TF, and
Definition~\ref{def:ievery} states the fact. The other 17 sites of P3 are
class (c).
\end{remark}

\subsection{P4: an upper bound of a length across a call growing stores (18 sites)}

\emph{Program.} \code{note(out, v)} pushes onto \code{out}. A loop with
test \code{i + 1 < a.flags.length} calls \code{note(out, i)} and reads
\code{a.starts[i + 1]}, where \code{a.flags} and \code{a.starts} are
\code{@stable}. The precondition states
$\len{a.\mathit{flags}}=\len{a.\mathit{starts}}$,
$a.\mathit{flags}\sep\mathit{out}$ and $a.\mathit{starts}\sep\mathit{out}$.
The two separations are part of the pattern. Without them the program is
unsafe: called with \code{out} the array \code{a.flags}, each push
lengthens \code{a.flags}, the test stays true, and the read passes the
end of \code{a.starts} \status{observed: the case
\code{invalid-bounds-pattern-grow-upper}, under Node}.

\emph{Fact.} $\len{a.\mathit{flags}}\le\len{a.\mathit{starts}}$ across
the call.

\emph{In TF.} Under the effect $\{\val:\any,\ \st:\grow\}$ the fact is
lost, since that effect lets every store grow, \code{a.flags} included
(TF §7, the full version of its Example 1). The callee's writes are
to one array its parameter names, so TF states for it the effect
$\{\val:\any,\ \mathit{out}.\kw{store}:\any\}$. With
$K=a.\mathit{flags}\sep\mathit{out}\wedge a.\mathit{starts}\sep\mathit{out}$
the fact survives by \textsc{(Sep)} at the rigid paths
\status{paper: \textsc{(Sep)}, \textsc{(Rigid)}}. Class (c), given the
callee's regional effect.

\emph{Sites.} 4 in chain 1 (\code{starts[i + 1]} under
\code{i + 1 < flags.length}), 10 in \code{measurement.ts} (the advances
list's exact length across measuring calls), 4 in
\code{prepareRichInline}. Whether each callee's store writes are at
paths of its parameters, so that a regional effect states them, is not
tested \status{conj.}.

\emph{Evidence.} \code{goodRegion}, the program above with the regional
effect and the separations at the paths \code{a.flags} and
\code{a.starts}, is refused, and \code{goodLower} (a lower bound) is
verified \status{observed: \code{t01\_grow\_equal.ts}}. The variant with
no separation stated is refused, and that refusal is correct. \code{slots}, the regional effect with the
arrays as parameters, is verified \status{observed:
\code{t01b\_region.ts}}; \code{paths} is refused and \code{slots} verified
\status{observed: \code{t01c\_lists.ts}}. TScert keeps the fact for
variables and not for rigid paths.

\begin{remark}[Classification of P4]
The analysis places P4 in class (b), since the fact is lost under the
effect \code{stores: "grow"}. In TF a regional effect states the callee
exactly, and under it the fact is kept. This note places P4 in class (c).
\end{remark}

\subsection{P6: a fact over value locations across a call changing values (12 sites)}

\emph{Program.} A loop keeps
$\every(\mathit{starts},\len{\mathit{starts}},0\le v<\kw{view}(\mathit{offsets}))$
across \code{touch(c)}, stated with the effect $\{\val:\any\}$
\status{observed: \code{t14\_every\_call.ts}, \code{badAny}}.

\emph{Where TF stops.} $\val$ contains the element cells of
\code{starts}, and no axiom keeps the fact under $\{\val:\any\}$. The
function \code{touch} writes \code{c.value}, so TF states for it the
effect $\{c.\mathit{value}:\any\}$, under which the fact is kept by
kinds, as for P3. The capability is needed by a callee that writes
locations no path from its parameters reaches: module cells, and fields
of objects it does not receive. Its only TF effect is $\{\val:\any\}$ or
$\top$.

\begin{definition}[Kind regions]\label{def:kinds}
$\fld(\rho,h)=F$, $\cel(\rho,h)=C$ and $\elt(\rho,h)=E$. So $\val$
denotes the union of the three. Kind regions may also appear in the
bound of an effect variable (Definition~\ref{def:effvar}).
\end{definition}
\noindent
Example: \code{let hits = 0; function touch(): void \{ hits++; \}} has the
effect $\{\cel:\any\}$. That effect respects the support
$\{\mathit{starts}[\ast]:\keep\}$ under $\kw{true}$, since $C$ is
disjoint from $S\cup E$, so the $\every$ survives the call
\status{conj.: this program is not run}.

\emph{The worked program.} At line 8, \code{measure} writes the cell of
\code{calls}. Its effect in TF is $\{\val:\any\}$, which loses the loop
invariant $\every(\mathit{lists}.\mathit{kinds},\len{\mathit{lists}.\mathit{kinds}},0\le v<3)$.
With Definition~\ref{def:kinds} its effect is $\{\cel:\any\}$. The
invariant's support is $\{\mathit{lists}.\mathit{kinds}:\keep,\ \mathit{lists}.\mathit{kinds}[\ast]:\keep\}$,
of kinds $F$, $S$ and $E$, so it survives. The lengths in the invariant
survive too, by the same argument.

\emph{Sites.} 5 in \code{prepareRichInline} (\code{offsets[analysis.starts[\ldots]]}
across \code{measureAnalysis}), 3 in \code{materializeRichInlineLineRange},
1 in \code{getFollowingSpaceTail}, 3 in \code{getHanKerningTrims}. The
facts read \code{analysis.starts} through a field; under an entry
$\fld:\any$ they also need $\rigid(\mathit{analysis}.\mathit{starts})$,
which \code{@stable} gives. The 3 sites of \code{getHanKerningTrims} read
a field of a fresh object \code{out} that the callees do not receive;
kind regions keep that fact when the callees write no field location
\status{conj.}. Which locations pretext's measuring callees write is not
checked \status{conj.}.

\emph{Evidence.} \code{badAny} is refused; \code{badInit}, whose callee
has the effect $\{\val:\init\}$, is verified, as \textsc{(Init)} predicts
\status{observed: \code{t14\_every\_call.ts}}.

\begin{remark}[Classification of P6]
The minimal program of the analysis is class (c) in TF, by the effect
$\{c.\mathit{value}:\any\}$. The pattern is class (b) for callees whose
written locations no TF region names, as in the module-cell program
above. This note keeps P6 in class (b) on the analysis's reading that
pretext's measuring callees are of that kind \status{conj.}.
\end{remark}

\section{Facts of a location or of a type for good}\label{sec:cap4}

By (T4) a fact that reads only immutable locations is invariant. TF's
immutable locations are \code{@stable} fields and typed-array lengths;
the patterns here need more.

\subsection{P15: a module constant array (4 sites)}

\emph{Program.} \code{const KINDS: readonly string[] = ["a", "b", "c", "d"];}
read at \code{KINDS[f \& 3]} \status{observed: \code{t06\_const\_table.ts}}.

\emph{Fact.} $\len{\mathit{KINDS}}=4$ at each entry of the function.

\emph{Where TF stops.} The store location of the literal has a contract
that admits every store of string cells. A transition that shortens it is
admissible, so the fact is not $\bx\top$. The fact holds because no code
writes that store, which TF cannot state without ownership.

\begin{definition}[Immutable store and immutable content]\label{def:imm}
A store location $s$ has an \emph{immutable store} when its contract
admits one store $\sigma$. It has \emph{immutable content} when, besides,
each $l\in\mathit{cells}(\sigma)$ has a contract that admits one value.
The cell of a module \code{const} binding has the contract that admits
the value it is initialised with. An array literal held in a module
\code{const} binding of a \code{readonly} array type is allocated with
immutable content.
\end{definition}
\noindent
Example: in the worked program, cell 0 holds $k$ for good, the store
location 1 holds $\langle 2,3,4\rangle$ for good, and the locations 2, 3
and 4 hold their strings for good. So $F_2=(\len{\mathit{KINDS}}=3)$ and
$\every(\mathit{KINDS},3,v\text{ is a string})$ are $\bx\top$, which is
what line 30 needs at every entry of \code{lastKind}
\status{proved on paper: Section~\ref{sec:sound}}.

\emph{Sites.} \code{prepareWithSegments} (142), \code{hasProperty} (234),
\code{needsComplexTextPath} (89, 90).

\emph{Evidence.} \code{kindName} is refused; \code{kindNameLocal}, the
same literal declared in the function, is verified \status{observed:
\code{t06\_const\_table.ts}}.

\subsection{P5: bounds held by the contents of generated tables (16 sites)}

\emph{Program.} A module cell \code{let rows: Uint8Array | null = null}
is set by \code{rows ??= decode()}, where \code{decode} copies the bytes
of a packed string into a new \code{Uint8Array}. The function reads
\code{t[t[state]!]!}. With the table cached in the cell, the read is
first blocked by P16; with the table decoded at each call, it is refused
at \code{t[t[state]]} \status{observed: the case
\code{bounds-pattern-generated-table-refused}}.

\emph{Fact.} $\every(t,\kw{view}(t),0\le v<\kw{view}(t))$ for the decoded
table $t$.

\emph{Where TF stops.} Two things are missing. First, no step establishes
the fact: the decoder copies bytes without testing them. A check in the
decoder establishes it, or an assumption that the generated data meets
its generator's contract does (Section~\ref{sec:notcovered}). Second, the
element cells of the table are value locations, written by the decoder
after allocation, and any later step may write them. Definition~\ref{def:imm}
fixes content at allocation, and a decoded table gets its content after
allocation.

\begin{definition}[Sealed content]\label{def:seal}
A \emph{seal} of an array $z$ at a step is the statement that, from that
step on, the store location of $z$ and every cell its store lists have
the contract that admits the contents they hold at that step.
\end{definition}
\noindent
Example: \code{decode} seals \code{out} at its return. From then on
$\every(\mathit{out},\kw{view}(\mathit{out}),v<\mathit{states})$, once
checked, is $\bx\top$. A seal changes a location's contract, which TF
excludes (TF §10, \emph{Fixed contracts}). Its soundness is
\status{conj.} (Section~\ref{sec:sound}).

\emph{Sites.} In \code{line-breaks.ts} (\code{unpackTable},
\code{getClass}, \code{unpackClassRuns}, \code{unpackStateRows},
\code{markRuleBoundaries}), 12 in all; 2 in \code{gecko-line-breaks.ts};
1 in \code{getWebKitGenericFamilies}; 1 in \code{getScriptFacts}. Each
index is an entry or the length of a decoded table.

\emph{Evidence.} Constant-stride rows verify \status{observed:
\code{t09\_arith.ts}, \code{goodStride}}. No program with a decoded table
is run.

\subsection{P16: a cache location initialised at its first use (4 sites)}

\emph{Program.} \code{let bmp: Uint16Array | null = null;} and
\code{const table = (bmp ??= new Uint16Array(65536)); return table[cp]!;}
under $0\le\mathit{cp}<65536$ \status{observed: \code{t07\_lazy\_cell.ts}}.

\emph{Fact.} $\mathit{bmp}\doteq\code{null}\vee\kw{view}(\mathit{bmp})=65536$
at every entry (Definition~\ref{def:cellpath}).

\emph{Where TF stops.} Within one invocation \textsc{(Init)} keeps facts
across \code{??=} (TF §8, \emph{Write-once cache}). Across invocations
the fact is a property of the cell's contents, and the contract
\code{Uint16Array | null} does not fix it.

\begin{definition}[Refined cell contract]\label{def:cellinv}
A module binding $b$ may state a \emph{refinement} $I_b(v)$, a formula
of one value $v$ that reads, besides $v$, only immutable locations. The
contract of the cell $l_b$ is the set of values of $b$'s type that
satisfy $I_b$. Each write to $b$ proves $I_b$ of the written value. Each
read of $b$ assumes it.
\end{definition}
\noindent
Example: $I_{\mathit{bmp}}(v)=(v\doteq\code{null}\vee\kw{view}(v)=65536)$.
The initial value \code{null} satisfies it, and the write in \code{??=}
writes a typed array of length 65536. The read \code{table[cp]} follows,
since a typed array's length is immutable. The same definition applies to
a field written once by \code{??=}, such as \code{measurement.kanaKerning}.

\emph{Sites.} \code{hasProperty} (232), \code{shouldBreakFast} (492),
\code{nextBreakablePosition} (722), \code{getKanaKerning} (955). For the
first three the table comes from \code{unpackTable}, so its length is
also a P5 fact.

\emph{Evidence.} \code{bitsAt} is refused; \code{bitsFresh}, the table
allocated where it is read, is verified \status{observed:
\code{t07\_lazy\_cell.ts}}.

\subsection{P20: a type invariant of a value read back from a Map (1 site)}

\emph{Program.} \code{type ClassTable = \{ readonly index: Uint16Array \}}
with the invariant \code{this.index.length === 4352}; a function returns
\code{cache.get(name)} of a \code{Map<string, ClassTable>}, or builds the
table and sets it \status{observed: \code{t18\_map\_cache.ts}}.

\emph{Fact.} $\kw{view}(t.\mathit{index})=4352$ for the value $t$ that
\code{cache.get(name)} returns.

\emph{Where TF stops.} TF has no invariant of a type. A value read from a
native's result carries its type and no further fact.

\begin{definition}[Type contract]\label{def:tcon}
A type $T$ may state an invariant $I_T(v)$, a formula of one value $v$
that reads only immutable locations. The contract $T$ is the set of
values of type $T$ that satisfy $I_T$. Each step that gives a value the
type $T$ proves $I_T$ of it. Each read of a location of contract $T$
assumes it.
\end{definition}
\noindent
Example: with the field \code{index} of \code{ClassTable} declared
\code{@stable}, $I_{\mathit{ClassTable}}(v)=(\kw{view}(v.\mathit{index})=4352)$
reads only immutable locations. \code{build} gives its object literal the
type \code{ClassTable} and proves the invariant there. \code{cache.get(name)}
returns a value of type \code{ClassTable} or \code{undefined}, by the
contract of \code{Map}, so it returns a value that meets the invariant or
\code{undefined}.

\emph{Sites.} \code{markRuleBoundaries} (361), reading
\code{overrides.categories[i]} below \code{overrides.chars.length} of
rules cached in a \code{Map}.

\emph{Evidence.} \code{tableOf} is refused at its return;
\code{classOfBuilt}, the record its builder returns, is verified
\status{observed: \code{t18\_map\_cache.ts}}.

\section{Richer terms}\label{sec:cap5}

TF's atoms are linear, quantify one array at a time, and name elements
only through $\every$. The five patterns here need five term forms. Each
is class (a).

\subsection{P9: a product of two variables in an index (7 sites)}

\emph{Program.} \code{rows[state * width]} under
\code{rows.length === states * width} and $0\le\mathit{state}<\mathit{states}$
\status{observed: \code{t09b\_product.ts}, refused at the source gate}.

\begin{definition}[Row table]\label{def:row}
A term may be a product $x\cdot y$ of two variables. The atom
$\rowt(z,n,w)$ holds when $n,w\in\mathbb{N}$ and $\len z=n\cdot w$ (for a
typed array, $\kw{view}(z)=n\cdot w$). Its rule \textsc{(Row)} is
\[
\rowt(z,n,w)\wedge 0\le r<n\wedge 0\le c<w\ \vdash\ 0\le r\cdot w+c<\len z .
\]
\end{definition}
\noindent
Example: with $\rowt(\mathit{rows},\mathit{states},\mathit{width})$,
\code{row = state * width} and $3+\mathit{category}<\mathit{width}$, the
read \code{rows[row + 3 + category]} is below the length. The support of
$\rowt(z,n,w)$ is that of $\len z$.

\emph{Sites.} \code{parseGraphemeRules} (46, 47, 48),
\code{findGraphemeEnds} (85), \code{markRuleBoundaries} (369, 373, 382).
Each also needs the state read from the table to be below
\code{states}, a P5 fact.

\subsection{P10: equal counts computed by two functions of one input (7 sites)}

\emph{Program.} \code{ends(text)} and \code{widthsOf(text)} each push one
entry per non-space code unit of \code{text}; the caller reads
\code{w[k]} for $k<\len e$ \status{observed: \code{t12\_uf.ts}}.

\emph{Fact.} $\len e=\len w$.

\begin{definition}[Named count]\label{def:count}
A \emph{count} $g$ is declared by a prefix recurrence on strings:
$g(x,0)=0$ and $g(x,i+1)=g(x,i)+\beta(x,i)$ for $i<\slen(x)$, where
$\beta(x,i)\in\{0,1\}$ is a condition on the code unit of $x$ at $i$. It
names $g(x)=g(x,\slen(x))$. Terms may contain $g(x)$ and $g(x,t)$ for a
variable $x$ and a linear term $t$.
\end{definition}
\noindent
Example: $\beta(x,i)=1$ when the code unit at $i$ is not 32. Both
builders prove the loop invariant $\len{\mathit{out}}=g(\mathit{text},i)$,
and so the post $\len r=g(\mathit{text})$; the caller derives
$\len e=g(\mathit{text})=\len w$. The support of $g(x)$ is $\emptyset$.

\emph{Sites.} 6 in \code{measureAnalysis}, among them
\code{fitAdvances[k]} under \code{k < extraSpacings.length}; 1 in
\code{prepareRichInline} (523). Here $g$ counts graphemes. Whether
pretext's grapheme count has the form of Definition~\ref{def:count} is
not checked \status{conj.}.

\subsection{P12: a pointwise relation of two arrays (5 sites)}

\emph{Program.} Under
\code{forall k in fits: extras[k] === null || extras[k].length === fits[k].length},
the read \code{e[g]} with \code{e = extras[i]} and
\code{g < fits[i].length} \status{observed: \code{t04\_pointwise.ts},
refused at the source gate}.

\begin{definition}[Indexed every]\label{def:ievery}
The atom $\every_k(t,\varphi)$ holds when, for each $k$ with $0\le k<t$,
$\varphi$ holds and every read in $\varphi$ is defined. The formula
$\varphi$ may read, for several arrays $x$, the element $x[k]$, the
length $\len{x[k]}$ of an element that is an array, a field $x[k].f$, and
the guard $x[k]\doteq\code{null}$. Two regions name what it reads:
$z[\ast].f$ denotes the locations $u(f)$ of the objects $u$ held by the
cells the store of $z$ lists, and $z[\ast].\kw{store}$ denotes the store
locations of the arrays held by those cells.
\end{definition}
\noindent
Example:
$\every_k(\len{\mathit{fits}},\ \mathit{extras}[k]\doteq\code{null}\vee\len{\mathit{extras}[k]}=\len{\mathit{fits}[k]})$
bounds \code{e[g]} by \code{g < fits[i].length}. Its support is
$\{\mathit{fits}[\ast]:\keep,\ \mathit{extras}[\ast]:\keep,\ \mathit{fits}[\ast].\kw{store}:\keep,\ \mathit{extras}[\ast].\kw{store}:\keep\}$.
TF's $\every(z,t,\varphi)$ is $\every_k(t,\varphi[z[k]/v])$ with one
array. The 2 sites of P3 at line 1185 need
$\every_k(t,\mathit{fragments}[k].\mathit{itemIndex}<n)$, of support
$\{\mathit{fragments}[\ast]:\keep,\ \mathit{fragments}[\ast].\mathit{itemIndex}:\keep\}$.

\emph{Sites.} \code{countPreparedLines} (376),
\code{walkPreparedComplexLines} (1027),
\code{stepPreparedSimpleLineGeometry} (1134), \code{getPartWidth} (1001),
\code{getWebKitParagraphBreaks} (704). The last relates
$\mathit{starts}[k+1]-\mathit{starts}[k]$ to
$\slen(\mathit{items}[k].\mathit{text})$.

\subsection{P17: an index into the element another index selects (3 sites)}

\emph{Program.} \code{fits[c.segment]![c.grapheme]!} under
\code{c.grapheme < fits[c.segment].length} \status{observed:
\code{t11\_cursor.ts}, refused at the source gate}.

\begin{definition}[Indexed path]\label{def:ipath}
Paths are $\pi::=x\mid\pi.f\mid\pi[t]$ with $t$ a linear term. The read
of $\pi[t]$ is the element at index $t$ of the array that $\pi$ reads; it
fails at a hole or past the end. The region of the read of $\pi[t]$ is
the store location of that array with the cell at index $t$.
\end{definition}
\noindent
Example: $c.\mathit{grapheme}<\len{\mathit{fits}[c.\mathit{segment}]}$
reads the fields of \code{c}, the store of \code{fits}, the cell at
$c.\mathit{segment}$, and the store of the array that cell holds. Its
support gives each of these $\keep$, but the last, which tolerates
$\grow$.

\emph{Sites.} \code{stepPreparedSimpleLineGeometry} (1134, 1134) and
\code{getPartWidth} (998).

\subsection{P19: a running sum of element lengths (2 sites)}

\emph{Program.} A loop
\code{for (let i = 0, at = 0; i < segments.length; at += segments[i++]!.length)}
reads \code{offsets[at]}, with
$\kw{view}(\mathit{offsets})$ the sum of the segments' lengths, the view
being allocated at that sum. The read is refused \status{observed: the
case \code{bounds-pattern-running-sum-refused}}.

\begin{definition}[Element sum]\label{def:sum}
For an array $z$, terms $i,j$ and a term $\tau$ of one element,
$\esum(z,i,j,\tau)=\sum_{i\le k<j}\tau(z[k])$. Linear atoms may contain
it. Its rules are $\esum(z,j,j,\tau)=0$ and
$\esum(z,i,j,\tau)=\tau(z[i])+\esum(z,i+1,j,\tau)$ for $i<j$.
\end{definition}
\noindent
Example: the loop invariant
$\mathit{at}+\esum(\mathit{segments},i,\len{\mathit{segments}},\slen)=\kw{view}(\mathit{offsets})$
with each element nonempty gives $\mathit{at}<\kw{view}(\mathit{offsets})$
while $i<\len{\mathit{segments}}$. Its support is
$\{\mathit{segments}[\ast]:\keep\}$: a string element is a value.

\emph{Sites.} \code{findSegmentSources} (1270, 1271).

\section{Rules TF has and TScert lacks}\label{sec:cap6}

Each pattern here is class (c): the fact is stated, and TF derives it,
from rules TF states or from posts sound for the step.

\subsection{P7: a fact over a local closure's captures (10 sites)}

\emph{Program.} \code{const glyphs = [\ldots 4 objects \ldots];} and a
closure \code{group(from, to)} that reads \code{glyphs[i]!.advance} for
$\mathit{from}\le i<\mathit{to}$, called as \code{group(0, 2) + group(2, 4)}
\status{observed: \code{t19\_closure.ts}}.

\emph{In TF.} \code{glyphs} is a \code{const} that no closure writes, so
it is a variable (TF §2). The closure's precondition
$\len{\mathit{glyphs}}=4\wedge0\le\mathit{from}\wedge\mathit{to}\le4$
names it free, \textsc{Call} substitutes the parameters only, and
\textsc{Frame} carries $\len{\mathit{glyphs}}=4$ from the closure's
creation to each call \status{paper: TF §2, \textsc{Call},
\textsc{Frame}}.

\emph{Sites.} \code{getHanKerningTrims} (212, 213), \code{getFontData}
(140 four times, 142), \code{prepareRichInline} (577, 586, 620).

\emph{Evidence.} \code{groups} is refused; \code{inline}, the same reads
in the function's own body, is verified \status{observed:
\code{t19\_closure.ts}}.

\subsection{P11: a string's element read below its length (5 sites)}

\emph{Program.} \code{for (let i = 0; i < text.length; i++) out = text[i]!;}
\status{observed: \code{t08\_strings.ts}}.

\emph{In TF.} The fact $0\le i<\slen(\mathit{text})$ is over variables and
a string length, invariant by Definition~\ref{def:slen}, and the loop
test establishes it. The step needs its post.

\begin{proposition}[String read]\label{prop:strread}\status{proved on paper}
$\hoare{0\le i<\slen(s)}{y=s[i]}{y\text{ is a string}\wedge\slen(y)=1}$.
\end{proposition}
\begin{proof}
For a string and an integer index below its length, JavaScript's indexed
read yields the one-unit string of the code unit at that index.
\end{proof}

\emph{Sites.} \code{segmentAtLineBreaks} (390, 396),
\code{getWebKitLineRules} (787), \code{measureKanaKerning} (963, 963).

\emph{Evidence.} \code{badElement} is refused; \code{goodCode}, a
\code{charCodeAt(i)} in the same loop, is verified \status{observed:
\code{t08\_strings.ts}}.


\subsection{P13: a typed-array prefix filled with a postfix counter (5 sites)}

\emph{Program.} \code{for (let i = 1; i <= n; i++) ends[count++] = i;}
for \code{ends: Int32Array} with $n\le\kw{view}(\mathit{ends})$, returning
\code{count}, with the prefix fact stated as the post. The loop's back
edge is refused: the fact over the parameter \code{ends} is not extended
across the write. The same loop over a view allocated in the function is
verified \status{observed: the cases
\code{bounds-pattern-postfix-prefix-refused} and
\code{bounds-pattern-postfix-prefix}}.

\emph{Fact.} $\every(\mathit{ends},\mathit{count},0<v\le\slen(\mathit{text}))$,
established by the writes and returned in a post that names
\code{count}.

\begin{proposition}[Extending an every]\label{prop:extend}\status{proved on paper}
Let $\varphi$ be a formula of one element whose other terms are
invariant, and $t$ a variable.
\begin{enumerate}[label=(\alph*),itemsep=1pt,topsep=2pt]
\item For a typed array $z$:
$\hoare{\every(z,t,\varphi)\wedge0\le t<\kw{view}(z)\wedge\varphi(v)}{z[t]=v}{\every(z,t+1,\varphi)}$.
\item For an array $z$:
$\hoare{\every(z,\len z,\varphi)\wedge\dense(z,\len z)\wedge\varphi(v)}{z.\code{push}(v)}{\every(z,\len z,\varphi)\wedge\dense(z,\len z)}$.
\end{enumerate}
\end{proposition}
\begin{proof}
(a) The write changes the one element cell at index $t$. The cells at the
indices below $t$ are other cells of the same store, so they keep their
contents, and the cell at $t$ holds $v$. (b) The push changes the store
of $z$ by one new cell at index $\old(\len z)$, holding $v$, and leaves
the other cells. In both cases $\varphi$ reads nothing else that changes.
\end{proof}

\noindent
With the renaming $\mathit{count}'=\mathit{count}+1$ of the postfix
increment, (a) and \textsc{Loop} keep the fact. In the worked program,
(b) at line 10 keeps
$\every(\mathit{lists}.\mathit{kinds},\len{\mathit{lists}.\mathit{kinds}},0\le v<3)$,
since the pushed value is 0, 1 or 2.

\emph{Sites.} \code{setupClusterBoundaries} (353),
\code{getUnspacedGraphemes} (263), \code{prepareRichInline} (698, 699,
700).

\begin{remark}[Establishing an every]
The analysis attributes P13 to \textsc{Loop}. \textsc{Loop} keeps the
fact; Proposition~\ref{prop:extend} establishes it at each write. TF
states neither part of that proposition. The pattern is class (c).
\end{remark}

\subsection{P18: a guarded fact through two fields (2 sites)}

\emph{Program.} A type \code{Prepared} with \code{@stable} fields
\code{widths} and \code{items: Items | null}, \code{Items} with a
\code{@stable} field \code{insideExtras: number[] | null}, and the
invariant
$g\vee\len{p.\mathit{items}.\mathit{insideExtras}}=\len{p.\mathit{widths}}$
with $g=(p.\mathit{items}\doteq\code{null}\vee p.\mathit{items}.\mathit{insideExtras}\doteq\code{null})$
\status{observed: \code{t05\_depth2\_guard.ts}}.

\emph{In TF.} $\rigid(p.\mathit{items})$ holds, and under
$\neg(p.\mathit{items}\doteq\code{null})$ so does
$\rigid(p.\mathit{items}.\mathit{insideExtras})$. So $g$ is invariant,
and law (d) of TF's Proposition \emph{Laws},
$g\vee\bx\Phi F\vdash\bx\Phi(g\vee F)$, moves the guard under the
modality. In \code{badNoCall} the loop body only reads, its effect is
$\emptyset$, and the fact is kept \status{paper: Proposition
\emph{Laws} (d), \textsc{(Rigid)}, \textsc{Loop}}.

\emph{Sites.} \code{stepPreparedSimpleLineGeometry} (1138, 1138).

\emph{Evidence.} \code{bad} and \code{badNoCall} are refused; \code{good},
a guard of one field, is verified \status{observed:
\code{t05\_depth2\_guard.ts}}.

\subsection{P21: shift left and bitwise or (0 sites as first blocker)}

\emph{Program.} \code{values[(hi <{}< 8) | (c \& 0xff)]} under
$0\le\mathit{hi}<4352$ and $\kw{view}(\mathit{values})=1114112$
\status{observed: \code{t09\_arith.ts}}.

\begin{proposition}[Shift and or]\label{prop:shift}\status{proved on paper}
For integers $0\le a<2^{23}$ and $0\le d\le 255$, \code{(a <{}< 8) | d}
equals $256a+d$.
\end{proposition}
\begin{proof}
$256a<2^{31}$, so the 32-bit shift is exact and its low 8 bits are zero.
The or with $d<256$ sets those bits to $d$.
\end{proof}

\noindent
With $d=c\mathbin{\&}255\in[0,255]$, the index is
$256\,\mathit{hi}+d<4352\cdot256=1114112$. The post is linear.

\emph{Sites.} None as first blocker; one as second blocker, in
\code{getClass} (138), first blocked by P5.

\emph{Evidence.} \code{badBits} is refused; \code{goodShift},
\code{index[c >{}> 8]}, is verified \status{observed:
\code{t09\_arith.ts}}.

\section{The classification}\label{sec:class}

The table gives, for each pattern, its sites, the analysis's class, the
class this note finds (Definition~\ref{def:classes}), and what it needs.

\medskip
\noindent
{\small
\begin{tabular}{@{}llrllp{0.36\textwidth}@{}}
\toprule
cap. & pattern & sites & analysis & found & needs\\
\midrule
1 & P1 & 50 & (b) & (b) & 47: Def.~\ref{def:compl}; 3: Def.~\ref{def:effvar}\\
2 & P2 & 31 & (c) & (c) & \textsc{Call} with a two-state post\\
2 & P8 & 10 & (c) & (c) & Prop.~\ref{prop:writelen}\\
2 & P14 & 4 & (c) & (c) & native posts; Def.~\ref{def:slen} for 1\\
3 & P3 & 19 & (c) & 17 (c), 2 (a) & Tolerant frame; 2: Def.~\ref{def:ievery}\\
3 & P4 & 18 & (b) & (c) & a regional callee effect, \textsc{(Sep)}, \textsc{(Rigid)}\\
3 & P6 & 12 & (b) & (b) & Def.~\ref{def:kinds}\\
4 & P5 & 16 & (d) & (d) & Def.~\ref{def:seal}, and a check or trust\\
4 & P15 & 4 & (d) & (d) & Def.~\ref{def:imm}\\
4 & P16 & 4 & (d) & (d) & Def.~\ref{def:cellinv}\\
4 & P20 & 1 & (d) & (d) & Def.~\ref{def:tcon}\\
5 & P9 & 7 & (a) & (a) & Def.~\ref{def:row}\\
5 & P10 & 7 & (a) & (a) & Def.~\ref{def:count}\\
5 & P12 & 5 & (a) & (a) & Def.~\ref{def:ievery}\\
5 & P17 & 3 & (a) & (a) & Def.~\ref{def:ipath}\\
5 & P19 & 2 & (a) & (a) & Def.~\ref{def:sum}\\
6 & P7 & 10 & (c) & (c) & TF §2, \textsc{Call}, \textsc{Frame}\\
6 & P11 & 5 & (c) & (c) & Prop.~\ref{prop:strread}, Def.~\ref{def:slen}\\
6 & P13 & 5 & (c) & (c) & Prop.~\ref{prop:extend}, \textsc{Loop}\\
6 & P18 & 2 & (c) & (c) & Laws (d), \textsc{(Rigid)}\\
6 & P21 & 0 & (c) & (c) & Prop.~\ref{prop:shift}\\
\bottomrule
\end{tabular}}
\medskip

\noindent
By class, the 215 blocked sites are: (a) 26, (b) 62, (c) 102, (d) 25.
The analysis's counts are (a) 24, (b) 80, (c) 86, (d) 25. The difference
is P4 (18 sites from (b) to (c)) and 2 sites of P3 (from (c) to (a)). The
patterns P11 and P14 are class (c) only with Definition~\ref{def:slen}.

\section{Soundness of the extensions}\label{sec:sound}

Each statement below says what makes one extension sound in TF's model:
which properties of admissible transitions it uses, and what must be
checked where. A statement marked \status{proved on paper} carries its
argument.

\begin{description}[leftmargin=1.5em,style=sameline]
\item[String length (Def.~\ref{def:slen})] \status{proved on paper} An
atom over variables and string lengths is invariant. Argument: strings
are scalars, so $\slen(x)$ depends on $\rho(x)$ alone; a step extends
$\rho$ and changes none of its bindings (TF §2). No (T$n$) is used.
\item[Complement and object regions (Def.~\ref{def:compl})]
\status{proved on paper} For a variable $x$, the effect
$\{\val\setminus x.\ast:\any,\ \st:\grow\}$ is rigid. Argument: by (T3)
$x.\ast$ denotes the same set in every heap. By admissibility a location
keeps its contract, hence its class. A location of $h$ changed by the
first of two steps and again by the second lies, in both, in
$\val\setminus x.\ast$ or in $\st$; transitivity follows as for TF's class
regions, using (T1). For a path $\pi$ that is not rigid, $\pi.\ast$ is
read in each step's own heap and the effect is not rigid in general.
\item[Bounded effect variable (Def.~\ref{def:effvar})]
\status{proved on paper} If a derivation of $f$'s
triple uses $\eff p$ only as the effect of the calls of $p$ and through
$\eff p\sqsubseteq\Psi$, then replacing $\eff p$ by a rigid
$\Psi'\sqsubseteq\Psi$ whose regions rebase to themselves gives a
derivation of the instance. Argument: by induction on the derivation.
\textsc{Frame} at a call of $p$ needs $\bx{\eff p}R$, which the
derivation obtained from $\bx\Psi R$ by law (c); law (c) gives
$\bx{\Psi'}R$ the same way. \textsc{Loop} and \textsc{Seq-Eff} need the
effect rigid and its regions rebased to the entry; $\Psi'$ is rigid, and
its regions are class regions, kind regions, or regions at rigid paths,
which rebase to themselves (TF §6, \emph{Rebasing}). With the argument
\code{undefined}, each call of $p$ is skipped by its guard
(\code{visit !== undefined} or \code{?.}) or throws before any change, so
the instance $\emptyset$ bounds it. It uses (T1). What must be
checked: at each call, $\Psi'\sqsubseteq\Psi$; for a visitor passed by
code outside the analysed sources, this is assumed.
\item[Kind regions (Def.~\ref{def:kinds})] \status{proved on paper}
$\fld$, $\cel$, $\elt$ and $\st$ denote pairwise disjoint sets, and an
effect over them is rigid. Argument: TF partitions $\mathrm{dom}\,h$
into $F\uplus C\uplus E\uplus S$, and a location's kind is fixed by what
allocated it: an own key, a binding, an array element or an array. The
rigidity argument of TF's class regions applies kind by kind. Respect by
kind needs $K=\kw{true}$. No (T$n$) beyond (T1) is used.
\item[Immutable store and content (Def.~\ref{def:imm})] Two statements.
\status{proved on paper} A fact that reads only locations with
single-value contracts is $\bx\top$: by (T4) each of those locations holds
the same contents in every admissible successor. \status{conj.} Every
step of the program respects the contracts given to a module literal:
the literal is allocated with its contents, so the contracts hold at
allocation; afterwards, the store or a cell is written only through a
reference of a mutable array type, and a \code{readonly} array type is
not assignable to a mutable one. This rests on the absence of casts and
of \code{any} in the sources, and on natives that receive the array not
writing it.
\item[Sealed content (Def.~\ref{def:seal})] \status{conj.} A seal is sound
when, at the sealing step, no reference to the array or its cells exists
at a mutable type outside the sealing function's result. A seal narrows a
location's contract after allocation. TF's worlds keep every contract
(TF §3), so a seal needs worlds whose extension may narrow a contract,
and a typing judgment that remains valid under narrowing. Neither is
given here. After a seal, the statement of Def.~\ref{def:imm} applies by
(T4).
\item[Refined cell contract (Def.~\ref{def:cellinv})]
\status{proved on paper} If every write to $b$ proves $I_b$ and the
initial value satisfies it, then $I_b(h(l_b))$ holds in every reachable
heap. Argument: $I_b$ reads besides $v$ only immutable locations, so by
(T4) its truth for a fixed $v$ does not change. The contract of $l_b$ is
fixed at allocation, and every admissible successor meets it (TF §3).
What must be checked: each write to $b$; all of them are in $b$'s module,
since an ES module's importers cannot assign its bindings.
\item[Type contract (Def.~\ref{def:tcon})] \status{proved on paper} If
each step that gives a value the type $T$ proves $I_T$, then every value
of contract $T$ in a reachable state meets $I_T$. Argument: as for the
refined cell contract, with (T3) for the field locations $I_T$ reads
through and (T4) for their contents. What must be checked: each
construction of a $T$, each assignment that types a value of another
type at $T$, and each native whose result has type $T$; the last is an
assumption about the native (here \code{Map.get}).
\item[Row table (Def.~\ref{def:row})] \status{proved on paper} Argument:
$r\cdot w+c\le(n-1)w+(w-1)=nw-1<\len z$. The atom reads what $\len z$
reads, so it inherits that support.
\item[Named count (Def.~\ref{def:count})] \status{proved on paper} $g$ is a
function of a string value, defined by recurrence on $i\le\slen(x)$, so
$g(x)$ reads no location and atoms over it are invariant. A post
$\len r=g(x)$ is proved by the builder like any linear post. What must be
checked: that the declared recurrence is the one each builder's loop
follows.
\item[Indexed every and element regions (Def.~\ref{def:ievery})]
\status{proved on paper} The stated support is a support. Argument: the
truth of $\every_k(t,\varphi)$ is determined by the stores of the arrays
$x$, the contents of the cells they list, and the field locations and
element stores those contents lead to; the support keeps each. The
regions $z[\ast].f$ and $z[\ast].\kw{store}$ are read in the heap a
change starts from, as TF's regions are. It uses (T3).
\item[Indexed path (Def.~\ref{def:ipath})] \status{proved on paper} A
fact through $\pi[t]$ is supported by $\keep$ on each field, store and
cell its read passes through, and by the change the fact tolerates on the
last store it reads. Argument: those locations determine the read. A
growth of a store may list a new cell at a present index: TF's $\preceq$
bounds the length, the dense prefixes and the newness of listed cells,
and does not fix the cell at an index. So a store whose cell the read
selects needs $\keep$.
\item[Element sum (Def.~\ref{def:sum})] \status{proved on paper} The two
rules are the definition of a finite sum. The support
$\{z[\ast]:\keep\}$ and the supports of $\tau$ determine every summand.
\item[Posts (Props.~\ref{prop:writelen}, \ref{prop:strread},
\ref{prop:extend}, \ref{prop:shift})] \status{proved on paper} Each is a
sound two-state post of its step, by the step's semantics, and enters a
proof as $Q_e$ of checked transfer (TF §7). Every other fact crosses the
step by \textsc{Frame}.
\end{description}

\section{What is not covered}\label{sec:notcovered}

\paragraph{The sites not placed.} 3 sites of $U$:
\begin{description}[leftmargin=1.5em,style=sameline]
\item[\code{analyzeText} (522, 524)] The reads \code{scanStarts[k]} below
\code{paragraph.atomic.length} and \code{sourceBreaks[after]} with
\code{after = scanStarts[k] + 1}. They need
$\len{\mathit{scanStarts}}=\len{\mathit{paragraph}.\mathit{atomic}}$ after
a callee pushes into \code{scanStarts}, and
$\every(\mathit{scanStarts},\len{\mathit{scanStarts}},v<\slen(r))$ for
the callee's result string $r$. Both are TF assertions with
Definition~\ref{def:slen}; the first is a P2 post when it needs $\old$.
Neither is tested.
\item[\code{readScriptRuns} (193)] The read \code{openBrackets[opened + 1]}
with \code{opened = openBrackets.lastIndexOf(x)}. It is in bounds because
the array holds pairs and \code{x} occurs at no odd index. The parity
$\exists m.\,\len{\mathit{openBrackets}}=2m$ is a TF assertion. The fact
at odd indices needs Definition~\ref{def:ievery} with a guard on $k$, and
the post of \code{lastIndexOf} \status{conj.}. Neither is tested.
\end{description}

\paragraph{Untested programs.} Every pattern's program has an observed
verdict, each read on the branch \code{single-pass} at \code{b7c8a559}.
The worked program's verdicts are in Section~\ref{sec:worked}. The
module-cell program of P6 is this note's and is not run. The 165 locally supported sites without a
verified shape are \status{conj.}. The assignment of each pretext site to
a pattern is \status{conj.}, except where the minimal program matches the
site's shape directly: P1, P2, P3, P11, P15, P16 and P18.

\paragraph{Trust decisions.} Three facts enter as assumptions, not as
proofs:
\begin{description}[leftmargin=1.5em,style=sameline]
\item[The bound on a visitor passed by the application] The walkers are
exported. The application passes \code{onLine}, and its effect is not
analysed. The bound $\Psi_{\mathit{onLine}}$ of Section~\ref{sec:cap1}
is an assumption about it.
\item[Native length contracts] The posts of \code{replace}, \code{split},
\code{map}, \code{subarray} and \code{slice} in Section~\ref{sec:cap2},
the post of \code{lastIndexOf}, and the type of \code{Map.get}'s result
in Section~\ref{sec:cap4}. Each is a contract of a native, audited and
assumed.
\item[Generated tables] The entries of the tables that pretext's
generator packs meet that generator's contract: a back reference, a run
length, a class or a state below its bound. The alternative to the
assumption is a check in each decoder.
\end{description}

\paragraph{Outside the account.} The 34 rewrite markers whose operand is
no index read (a nullable binding, a field, a call result), and the
rewritten source's other changes, are outside this note.

\end{document}
